\section{Approach}
\label{sec:approach}

\subsection{Graph Construction}

\paragraph{Compile-time (CT) graph.}
We parse eleven \sqlite{} C files with \clang~18.1.3 via libclang, recording
every direct function call with its source location, argument identifiers, and
return-value assignments.
The resulting graph \textsc{CT-v1} has 417 defined functions, 392 external (boundary) functions,
94 struct/union types, 954 fields, and 2{,}438 call-site edges (1{,}593 distinct caller--callee pairs, including calls into external functions) stored in a SQLite database using
a shared schema (\texttt{nodes} and \texttt{edges} tables).
Forty indirect (function-pointer) calls are recorded in node metadata but not as edges.
A lightweight taint pass propagates three source patterns
(\texttt{sqlite3\_prepare*}, \texttt{sqlite3\_bind\_text*}, \texttt{sqlite3\_bind\_blob*}) from 11 source functions
to the sink functions in three classes (copy, alloc, free); of the 16 sinks defined, 6 are reached.
Depth-limited search ($K=6$) adds 84 \texttt{taint\_reaches} edge records (41 distinct source--sink pairs) to produce \textsc{CT-v2}.
Struct and union fields are stored separately as \texttt{has\_field\_of\_type} edges for type-hierarchy queries.
The taint pass is a purely syntactic over-approximation: it does not model aliasing, pointer arithmetic, or global state.
Taint edges were not hand-validated on the final graph and no result relies on them.

\paragraph{Syntactic (SYN) graph.}
We run the \texttt{graphify} tool (tree-sitter back-end) on the same eleven
files and convert the resulting JSON to the shared schema via an adapter script.
The resulting graph contains 440 functions and 772 call edges (of which 310 were
inferred by heuristic).
Unlike \clang which models each distinct call site, \texttt{graphify} records unique
caller--callee pairs.
Furthermore, \texttt{graphify} failed to parse macro-heavy header definitions,
so it did not extract struct fields or taint information; thus the SYN graph
does not support type-hierarchy or taint queries.
End-line values missing from the tree-sitter output are recovered via a
brace-matching or next-definition heuristic.

\subsection{Tool Set}

The two graphs expose different tool subsets (Table~\ref{tab:tools}).
CT supports all seven tools. SYN supports five, because it has no type or taint edges.
Results are sorted and capped at 50 items; source output is capped at 8{,}000 characters (roughly 100--200 lines) per call, with paging for longer functions.
All tools return \texttt{\{"result":~\ldots,~"truncated":~bool\}} or
\texttt{\{"error":~CODE,~\ldots\}} and are exposed to the LLM as OpenAI
function-calling schemas.

\begin{table}[h]
\caption{Graph-navigation tools and their availability per graph.}
\label{tab:tools}
\small
\begin{tabular}{lp{4.0cm}cc}
\toprule
\textbf{Tool} & \textbf{Description} & \textbf{CT} & \textbf{SYN} \\
\midrule
\texttt{get\_callers(fn)}           & Direct callers of a function        & \checkmark & \checkmark \\
\texttt{get\_callees(fn)}           & Direct callees of a function        & \checkmark & \checkmark \\
\texttt{find\_call\_path(src, dst)} & BFS shortest call chain             & \checkmark & \checkmark \\
\texttt{get\_function\_source(fn)}  & Paginated source lines              & \checkmark & \checkmark \\
\texttt{search\_functions(substr)}  & Fuzzy name search                   & \checkmark & \checkmark \\
\texttt{get\_type\_hierarchy(type)} & Struct/union fields and types       & \checkmark & --- \\
\texttt{get\_taint\_paths(param)}   & Sinks reachable from tainted param  & \checkmark & --- \\
\bottomrule
\end{tabular}
\end{table}

\subsection{Agent Loop (TOOL Conditions)}

The agent runs up to 15 turns via the OpenAI-compatible function-calling
interface at temperature~0.
On each turn the model either calls a tool (observation appended to history) or
emits a final JSON answer conforming to the \textsc{Answer} schema (required
fields: \texttt{answer\_function}, \texttt{top3}, \texttt{evidence\_path},
\texttt{taint\_path}, \texttt{cited\_edges}, \texttt{explanation}).
If no final answer is produced within 15 turns, a force-prompt is appended
and one additional call is made without tools (\texttt{stopped\_by =
forced\_final}).
All 48 tool-condition runs---36 from \qwen{} (12 non-ablation, 24 ablation)
and 12 from \nemotron{}---hit the 15-turn cap; every reported tool-condition
answer is therefore a forced final answer.

\subsection{Text Slice (TEXT Conditions)}

A breadth-first search from the entry function over \texttt{calls} edges serialises the call neighbourhood as text, one block per function, up to $B=6{,}000$ tokens (tiktoken \texttt{cl100k\_base}).
Blocks are added whole and the slice stops before the first one that would exceed the ceiling.
Five questions share \texttt{sqlite3Select} as entry, so they share one slice (2{,}849 tokens on SYN, 4{,}405 on CT, both complete).
Q5 uses \texttt{sqlite3\_step} (2{,}963 tokens on SYN, 5{,}982 tokens on CT, cut by the ceiling).
The slice is prepended to the user message, before the question text.
The model answers in one turn without tool access.

\subsection{Evaluation Design}

Table~\ref{tab:conditions} summarises the nine conditions.
\textsc{NOCTX} provides neither code nor tools and is repeated three times
per question at temperature~0.7 to estimate baseline variance.
All other conditions use temperature~0 and run once per question.
Ablation conditions (last four rows) were run on \qwen only: the NVIDIA API
returned malformed or empty completions for \nemotron on every tool-subset
call; this problem did not manifest in the full tool-set conditions.

\begin{table}[h]
\caption{Experimental conditions. ``5 tools'' denotes the SYN-compatible subset;
  ``7 tools'' includes type-hierarchy and taint paths (CT only);
  ``src+search'' denotes \texttt{get\_function\_source} and
  \texttt{search\_functions}.}
\label{tab:conditions}
\small
\begin{tabular}{lllp{3.5cm}}
\toprule
\textbf{Condition} & \textbf{Graph} & \textbf{Delivery} & \textbf{Tools} \\
\midrule
\textsc{NOCTX}           & ---      & None       & None \\
\textsc{SYN-TOOL}         & SYN      & Tool calls & 5 tools \\
\textsc{CT-TOOL}          & CT-v2    & Tool calls & 7 tools \\
\textsc{SYN-TEXT}         & SYN      & Text slice & None \\
\textsc{CT-TEXT}          & CT-v2    & Text slice & None \\
\textsc{CT-ct\_calls}     & CT-v2    & Tool calls & 4 (callers, callees, source, search) \\
\textsc{CT-ct\_types}     & CT-v2    & Tool calls & 5 (+ type hierarchy) \\
\textsc{CT-src-only}      & CT-v2    & Tool calls & 2 (src+search) \\
\textsc{SYN-src-only}     & SYN      & Tool calls & 2 (src+search) \\
\bottomrule
\end{tabular}
\end{table}
